\subsection{Stability of normalized kernels}
\label{sec:normalized-kernel-stability}

The total kernel mass is the dimension of the selected space. Dividing by that dimension removes its growth from the parameter estimates. The squared reproducing kernel supplies a symmetric probability coupling whose positivity improves the resulting variation bound.

Keep the fixed data of \cref{sec:partial-kernel-parameters}, and assume $N_m>0$. For each real $v\in C(X)$ put
\begin{equation}\label{eq:normalized-kernel-probability}
 \mu_v=\frac{B_v}{N_m}\nu.
\end{equation}
For a weighted orthonormal basis $s_1,\ldots,s_{N_m}$ use the local unit evaluations $\ell_x$ to set
\[
 F_i(x)=\mathrm{e}^{-mv(x)/2}\ell_x(s_i),\qquad
 C_v(x,y)=\sum_i F_i(x)\overline{F_i(y)}.
\]
Although $C_v$ depends on the evaluation phases, its squared modulus does not. Off $Z_m$ let $\sigma_{v,x}$ be the probability in \eqref{eq:partial-kernel-response-probability} at weight $v$, and define
\[
 \mathcal Q_v f(x)=\int_X f(y)\,\mathrm{d}\sigma_{v,x}(y).
\]
Set this expression to zero on $Z_m$, which is $\mu_v$-null.

\begin{lemma}\label{lem:coherent-kernel-markov}\index{Bergman kernel!coherent-state coupling}
The measure
\begin{equation}\label{eq:coherent-kernel-coupling}
 \mathrm{d}J_v(x,y)=\frac{|C_v(x,y)|^2}{N_m}
                         \,\mathrm{d}\nu(x)\,\mathrm{d}\nu(y)
\end{equation}
is a symmetric probability coupling of $\mu_v$ with itself. The operator $\mathcal Q_v$ is a selfadjoint positive semidefinite contraction on $L^2(\mu_v)$, preserves constants almost everywhere, and preserves the integral.
\end{lemma}
\begin{proof}
Orthonormality gives $\int_X|C_v(x,y)|^2\,\mathrm{d}\nu(y)=B_v(x)$. Symmetry gives the other marginal, and the mass identity gives total mass one. Dividing by $B_v(x)$ off $Z_m$ gives the squared Riesz-vector density in \eqref{eq:partial-kernel-response-probability}. Thus the conditional probabilities of $J_v$ define $\mathcal Q_v$. Their mass and the invariant marginal prove the two preservation assertions; Jensen's inequality gives the $L^2$ contraction. Symmetry gives selfadjointness. For $f\in L^2(\mu_v)$, finite expansion gives
\[
 \langle\mathcal Q_v f,f\rangle_{L^2(\mu_v)}
 =\frac1{N_m}\sum_{i,j}
       \left|\int_X fF_i\overline{F_j}\,\mathrm{d}\nu\right|^2\geq0.
\]
The matrix features are phase independent. Their integrals exist because $|F_i\overline{F_j}|\leq B_v$ and $f\in L^1(\mu_v)$. This proves positivity on the stated domain.
\end{proof}

\Needspace{15\baselineskip}
\begin{proposition}\label{prop:normalized-kernel-weight-stability}\index{Bergman kernel!weight stability}
For $q,v,w\in C(X,\mathbb R)$, write $\operatorname{osc}_Xq=\sup_Xq-\inf_Xq$. Then
\begin{equation}\label{eq:normalized-kernel-response-bound}
 D\mu_v[q]=m(\mathcal Q_vq-q)\mu_v,
 \qquad
 \|D\mu_v[q]\|_{\mathrm{TV}}
       \leq\frac m2\operatorname{osc}_Xq.
\end{equation}
Consequently
\begin{equation}\label{eq:normalized-kernel-weight-stability}
 \|\mu_w-\mu_v\|_{\mathrm{TV}}
 \leq2\tanh\left(\frac m4\operatorname{osc}_X(w-v)\right).
\end{equation}
These bounds are independent of $N_m$. For a sequence of nonempty selected spaces, $m\operatorname{osc}_X(w_m-v_m)\to0$ makes the two probability sequences have the same weak or total-variation limit whenever either converges in that topology.
\end{proposition}
\begin{proof}
The signed measure identity follows from \cref{cor:partial-kernel-pointwise-response} and differentiation in total variation. Put $g=q-\mathcal Q_vq$, $A=\{g>0\}$ and $p=\mu_v(A)$. The mean of $g$ is zero. Selfadjointness gives
\[
 \int_Ag\,\mathrm{d}\mu_v
   =\langle q,(I-\mathcal Q_v)\mathbf1_A\rangle_{L^2(\mu_v)}.
\]
The function $(I-\mathcal Q_v)\mathbf1_A$ is nonnegative on $A$ and nonpositive on its complement, almost everywhere. Its integral is zero. Subtracting $\inf_Xq$ in this pairing therefore bounds it above by
\[
 \operatorname{osc}_Xq
    \bigl(p-\langle\mathbf1_A,\mathcal Q_v\mathbf1_A\rangle\bigr).
\]
Positivity, applied to $\mathbf1_A-p$, gives
$\langle\mathbf1_A,\mathcal Q_v\mathbf1_A\rangle\geq p^2$; the cross terms vanish by preservation of constants and the integral. Hence
\[
 \|g\|_{L^1(\mu_v)}
 =2\int_Ag\,\mathrm{d}\mu_v
 \leq2\operatorname{osc}_Xq\,p(1-p)
 \leq\tfrac12\operatorname{osc}_Xq.
\]
Here total variation is the full signed-measure norm, the integral of the absolute density. This proves \eqref{eq:normalized-kernel-response-bound}.

For the global estimate take $q=w-v$, $v_t=v+tq$, $\mu_t=\mu_{v_t}$ and $\mathcal Q_t=\mathcal Q_{v_t}$ for $0\leq t\leq1$. Put $a=m\operatorname{osc}_Xq$. For any Borel set $E$, its probability $p_t=\mu_t(E)$ is differentiable, since integration against $\mathbf1_E$ is a bounded linear functional on measures in total variation. Selfadjointness gives
\[
 p'_t=-m\langle q,(I-\mathcal Q_t)\mathbf1_E\rangle_{L^2(\mu_t)}.
\]
The function $(I-\mathcal Q_t)\mathbf1_E$ has positive and negative parts whose integrals both equal $p_t-\langle\mathbf1_E,\mathcal Q_t\mathbf1_E\rangle$. Bounding $q$ between its infimum and supremum on those pieces, and applying positivity to $\mathbf1_E-p_t$, gives
\[
 |p'_t|\leq a\bigl(p_t-\langle\mathbf1_E,\mathcal Q_t\mathbf1_E\rangle\bigr)
          \leq ap_t(1-p_t).
\]
Every kernel is positive precisely off the same common zero locus $Z_m$, so the probabilities $\mu_t$ are mutually absolutely continuous. If $p_0$ is zero or one, $p_t=p_0$ throughout. Otherwise $0<p_t<1$, and integration gives
\[
 \left|\log\frac{p_1}{1-p_1}-\log\frac{p_0}{1-p_0}\right|\leq a.
\]
With $\sigma(z)=\mathrm{e}^z/(1+\mathrm{e}^z)$, if $p_1\geq p_0$ and $z=\log(p_0/(1-p_0))$, then
\[
 p_1-p_0\leq\sigma(z+a)-\sigma(z)
 =\frac{\sinh(a/2)}{\cosh(z+a/2)+\cosh(a/2)}
 \leq\tanh(a/4).
\]
Swapping the endpoints handles the other order. The zero-mass signed difference has norm $2\sup_E|\mu_w(E)-\mu_v(E)|$, proving \eqref{eq:normalized-kernel-weight-stability}. Its right-hand side tends to zero under the sequence hypothesis, giving the final assertion.
\end{proof}

\Needspace{12\baselineskip}
\begin{example}\label{ex:sharp-kernel-weight-modulus}\index{Bergman kernel!optimal oscillation modulus}
For each $a>0$ and $m\geq1$, the bound \eqref{eq:normalized-kernel-weight-stability} is optimal as a supremum over smooth weights, already when the selected space has dimension one. On $X=\mathbb P^1$, take the trivial line bundle with unit metric, selecting potential zero and normalized Fubini--Study area $\nu$. There are smooth real weights $w_j$ with $m\operatorname{osc}_Xw_j=a$ such that
\[
 \|\mu_{w_j}-\mu_0\|_{\mathrm{TV}}\longrightarrow2\tanh(a/4).
\]
\end{example}
\begin{proof}
The selected space consists of constants, so its probability at weight $u$ is $\mathrm{e}^{-mu}\nu/\int_X\mathrm{e}^{-mu}\,\mathrm{d}\nu$. Choose a spherical cap $E$ of mass $p=(1+\mathrm{e}^{a/2})^{-1}$. Its smooth boundary has zero area. Choose smooth cutoffs $0\leq\chi_j\leq1$, equal to one on an inner cap and zero outside an outer cap, converging almost everywhere to $\mathbf1_E$. Their oscillations are one. Put $w_j=-a\chi_j/m$. Bounded convergence of their exponential densities and positive normalizing constants gives total-variation convergence to
\[
 \frac{\mathrm{e}^{a\mathbf1_E}}{Z}\nu,
 \qquad Z=1-p+\mathrm{e}^ap=\mathrm{e}^{a/2}.
\]
The limiting mass of $E$ is $1-p$. Its density exceeds one on $E$ and is less than one outside. Its distance from $\nu$ is therefore $2(1-2p)=2\tanh(a/4)$. The triangle inequality transfers this limit to the smooth weights. This proves supremal sharpness without an assertion of exact smooth attainment.
\end{proof}

At the inverse-degree scale, the analytic control also becomes independent of the selected-space dimension.

\Needspace{13\baselineskip}
\begin{proposition}\label{prop:normalized-kernel-uniform-radius}
Put $\rho=\log(3/2)$. The map $f\mapsto\mu_{v+f/m}$ has a holomorphic continuation to $\{f\in C(X,\mathbb C):\|f\|_\infty<\rho\}$ with total-variation norm at most three. For real $q\in C(X)$ and $k\geq1$,
\begin{equation}\label{eq:normalized-kernel-cauchy-bounds}
 \left\|\left.\frac{\mathrm{d}^k}{\mathrm{d}t^k}\right|_{t=0}
                      \mu_{v+tq}\right\|_{\mathrm{TV}}
 \leq3k!\left(\frac{m\operatorname{osc}_Xq}{2\rho}\right)^k.
\end{equation}
The radius and norm bound are independent of $m,N_m$ and the real base weight $v$.
\end{proposition}
\begin{proof}
On $V_m$ use $T(f)=P_v(\mathrm{e}^{-f}\,\cdot\,)$. The proof of \cref{prop:partial-kernel-weight-analyticity} gives $\|T(f)^{-1}\|\leq2$ on the stated ball. Its continued measure density is
\[
 \frac{\mathrm{e}^{-mv(x)-f(x)}}{N_m}
         \ell_x\bigl(T(f)^{-1}k_{v,x}\bigr).
\]
The fixed Riesz identity bounds its absolute value by
$2\mathrm{e}^{\|f\|_\infty}B_v(x)/N_m$. Integration bounds the total variation by three; holomorphy follows from the same finite-space construction and linear passage to measures. For any complex constant $c$, $T(f+c)=\mathrm{e}^{-c}T(f)$, and the opposite factor in the density cancels. For nonconstant real $q$, replace it by $q-c$ with $c=(\sup_Xq+\inf_Xq)/2$, whose norm is $\operatorname{osc}_Xq/2$. Apply the Banach-valued Cauchy estimate on every circle of radius less than $\rho/(m\|q-c\|_\infty)$, then let the radii increase to that value. This proves \eqref{eq:normalized-kernel-cauchy-bounds}. A constant $q$ leaves the measure fixed and gives zero derivatives.
\end{proof}

The condition in \eqref{eq:normalized-kernel-weight-stability} cannot in general be replaced by uniform convergence of the weights. A fiber norm can retain an inverse-degree perturbation in a nonzero measure limit.

\Needspace{15\baselineskip}
\begin{example}\label{ex:inverse-degree-fiber-weight}\index{Bergman kernel!inverse-degree weight perturbation}
Let $Y$ be a connected compact complex manifold of positive dimension with a smooth positive probability volume $\nu_Y$. On $X=\mathbb P^1\times Y$, pull back $\mathcal O(1)$ with its Fubini--Study metric, take the selecting potential zero, and use $\nu=\nu_{\mathrm{FS}}\otimes\nu_Y$, where $\nu_{\mathrm{FS}}$ is the Fubini--Study area of mass one. Choose a nonconstant $q\in C(Y,\mathbb R)$ and put
\[
 v_m(b,y)=q(y)/m,\qquad
 Z_q=\int_Y\mathrm{e}^{-q}\,\mathrm{d}\nu_Y,\qquad
 \nu_q=Z_q^{-1}\mathrm{e}^{-q}\nu_Y.
\]
Then $N_m=m+1$ and, exactly for every $m$,
\begin{equation}\label{eq:inverse-degree-fiber-kernels}
 B_{m,v_m}(b,y)=(m+1)\frac{\mathrm{e}^{-q(y)}}{Z_q},
 \qquad
 \mu_{m,v_m}=\nu_{\mathrm{FS}}\otimes\nu_q.
\end{equation}
Thus $v_m\to0$ uniformly, but the $r=1$ kernel measures tend in total variation to $\nu_{\mathrm{FS}}\otimes\nu_q$, whereas the limit at weight zero is $\nu_{\mathrm{FS}}\otimes\nu_Y$. Both logarithmic kernel potentials tend uniformly to zero. Normalization in ambient dimension gives zero in total variation.
\end{example}
\begin{proof}
Every section descends by compact connected fibers, as in \cref{prop:partial-kernel-submersion}. The base space $H^0(\mathbb P^1,\mathcal O(m))$ has dimension $m+1$. Its unweighted kernel is constant by the transitive unitary action preserving the metric and volume; its integral gives that constant. The product weighted norm is $Z_q$ times the base norm, proving \eqref{eq:inverse-degree-fiber-kernels}. At zero weight the probability is $\nu_{\mathrm{FS}}\otimes\nu_Y$. These probabilities differ: equality would force $\mathrm{e}^{-q}/Z_q=1$ almost everywhere and then everywhere, since every nonempty open set has positive volume and $q$ is continuous, contrary to the choice of $q$. The $r=1$ measures are $(m+1)/m$ times their probabilities. Ambient normalization in dimension $1+\dim Y$ tends to zero in total variation. Finally the difference of logarithmic kernel potentials is $(-q-\log Z_q)/m$, which tends uniformly to zero.
\end{proof}

This fiber effect is the weight form of the retained conditional probabilities in \cref{sec:horizontal-kernel-asymptotics}. The stability estimates varied the norm within one selected space. A modification also changes the integrability data, so the next comparison must transform both the space and its norm.
